\documentclass[10pt,a4paper]{amsart}
\usepackage[margin=19mm]{geometry}
\usepackage{amsmath,amssymb,mathtools}
\usepackage[scr=rsfs,cal=euler]{mathalfa}
\usepackage{xfrac}
\usepackage[unicode,hidelinks,bookmarks=false]{hyperref}
\newcommand{\ie}{i.e.\ }
\newcommand{\cf}{cf.\ }
\newcommand{\resp}{respectively\ }
\newcommand{\Subsection}{\subsection}
\providecommand{\nobreakdash}{\nobreak\hbox{-}}
\setlength{\emergencystretch}{2em}
\multlinegap0pt
\usepackage{siunitx}
\usepackage{siunitx}
\usepackage{siunitx}
\usepackage{siunitx}
\usepackage{xcolor}
\usepackage{algorithm}\usepackage{algpseudocode}
\usepackage{tcolorbox}
\usepackage[shortlabels]{enumitem}
\usepackage{float}
\usepackage{siunitx}
\usepackage{amssymb}
\usepackage{mathtools}
\usepackage{tikz}
\newtheorem{lem}{Lemma}
\newcommand{\Norm}[1]{\left\| #1 \right\|}
\newcommand{\abs}[1]{\left| #1 \right|}
\providecommand{\norm}[1]{\Norm{#1}}
\newcommand{\reals}{\mathbb{R}}
\title{The Fast Mixing Mechanism for Differential Privacy}
\author{Omri Lev, Moshe Shenfeld, Vishwak Srinivasan, Katrina Ligett, Ashia C. Wilson}
\date{Source: arXiv:2605.30600v2 (CC BY 4.0)}
\begin{document}
\maketitle

\noindent\emph{Source and licence.} The Fast Mixing Mechanism for Differential Privacy. Version arXiv:2605.30600v2 (CC BY 4.0), distributed by arXiv under the Creative Commons Attribution 4.0 International licence, http://creativecommons.org/licenses/by/4.0/, as recorded in the arXiv licence metadata for this exact version and shown on the abstract page https://arxiv.org/abs/2605.30600v2. Full licence notice: \texttt{licenses/arxiv-2605.30600-CC-BY-4.0.txt}.\par\smallskip

\noindent\textit{Editorial scope.} Counted unit: the article's lem lem:operator\_norms (source0.tex lines 1996--2005). The article's complete original proof of it is delivered with it (source0.tex lines 2007--2108, one \texttt{proof} environment of 102 lines). No mathematics is written, repaired or re-derived: the statement and the proof are the article's own lines, in the article's own notation, and no retained line is substituted or re-flowed.  No further source-local context is needed: every cross-reference of every retained line resolves inside the two delivered spans.  Nothing was omitted from any retained span, so the package carries no omission record.  No retained line contains a citation, so the wrapper carries no bibliography and no bibliography entry.  No retained line contains a figure environment, an image-inclusion command or drawing code, so no image is delivered and the package declares no asset dependency.  Editorial formatting only, in addition: an AMS \texttt{amsart} preamble that restates the article's own declarations and macros used by the retained lines, namely the environments \texttt{lem} on separate counters, together with the standard packages those lines need.  The retained environments print in this document as Lemma 1 in order of appearance, whereas the article prints the same items with the numbering of its own full document.  The complete native source of the article as distributed in the arXiv e-print for this exact version is bundled unchanged as \texttt{sources/201682/source0.tex}.
\begin{lem}
    \label{lem:operator_norms}
    Let $u,v\in\reals^{d}$. Then,
    \begin{align}
        \norm{uv^{\top}}_{\mathrm{op}} = \norm{u}\norm{v},
        \qquad
        \norm{uv^{\top} + vu^{\top}}_{\mathrm{op}}
        = \norm{u}\norm{v} + \abs{u^{\top}v}.
    \end{align}
\end{lem}

\begin{proof}
    If either $u=\vec{0}_d$ or $v=\vec{0}_d$, both identities are immediate. Hence, assume
    $u,v\neq \vec{0}_d$, and define
    \[
        \overline{u}\coloneq \frac{u}{\norm{u}},
        \qquad
        \overline{v}\coloneq \frac{v}{\norm{v}}.
    \]

    We first calculate the operator norm of $uv^{\top}$. We have
    \begin{align}
        uv^{\top}
        =
        \norm{u}\norm{v}\,
        \overline{u}\ \overline{v}^{\top}.
    \end{align}
    Since $\overline{u}\ \overline{v}^{\top}$ is a rank-one matrix with its only
    nonzero singular value equal to one, it follows that
    \begin{align}
        \norm{uv^{\top}}_{\mathrm{op}}
        =
        \norm{u}\norm{v}.
    \end{align}

    Next, consider
    \begin{align}
        A
        \coloneq
        uv^{\top}+vu^{\top}
        =
        \norm{u}\norm{v}
        \left(
        \overline{u}\ \overline{v}^{\top}
        +
        \overline{v}\ \overline{u}^{\top}
        \right).
    \end{align}
    Then, 
    \begin{align}
        \norm{A}_{\mathrm{op}} = \max_{\norm{x}=1}\abs{x^{\top}Ax} = 2\norm{u}\norm{v}\max_{\norm{x}=1}\abs{(\overline{u}^{\top}x)(\overline{v}^{\top}x)}.
    \end{align}
    Thus, it remains to show that
    \begin{align}
        \max_{\norm{x}=1}
        \abs{
        (\overline{u}^{\top}x)(\overline{v}^{\top}x)
        }
        =
        \frac{1+\abs{\overline{u}^{\top}\overline{v}}}{2}.
    \end{align}

    Let $\rho\coloneq \abs{\overline{u}^{\top}\overline{v}}$. For any
    $\norm{x}=1$, by $2|ab|\leq a^{2}+b^{2}$,
    \begin{align}
        2\abs{
        (\overline{u}^{\top}x)(\overline{v}^{\top}x)
        }
        &\leq
        (\overline{u}^{\top}x)^{2}
        +
        (\overline{v}^{\top}x)^{2} \\
        &=
        x^{\top}
        \left(
        \overline{u}\ \overline{u}^{\top}
        +
        \overline{v}\ \overline{v}^{\top}
        \right)x \\
        &\leq
        \lambda_{\max}
        \left(
        \overline{u}\ \overline{u}^{\top}
        +
        \overline{v}\ \overline{v}^{\top}
        \right).
    \end{align}
    The nonzero eigenvalues of
    $\overline{u}\ \overline{u}^{\top}
    +
    \overline{v}\ \overline{v}^{\top}$
    are $1+\rho$ and $1-\rho$, which follows since the eigenvectors of this matrix are $\frac{\overline{u} + \overline{v}}{\norm{\overline{u} + \overline{v}}}$ and $\frac{\overline{u} - \overline{v}}{\norm{\overline{u} - \overline{v}}}$. Hence,
    \begin{align}
        \max_{\norm{x}=1}
        \abs{
        (\overline{u}^{\top}x)(\overline{v}^{\top}x)
        }
        =
        \frac{1+\rho}{2}.
    \end{align}
    Substituting this into the previous equation gives
    \begin{align}
        \norm{uv^{\top}+vu^{\top}}_{\mathrm{op}}
        &=
        2\norm{u}\norm{v}\cdot
        \frac{1+\abs{\overline{u}^{\top}\overline{v}}}{2} \\
        &=
        \norm{u}\norm{v}
        +
        \abs{u^{\top}v}.
    \end{align}
    This completes the proof.
\end{proof}

\end{document}
